\chapter{CONCLUSION AND FUTURE WORK}
\label{chap:conclusion}

Despite of the final thesis, which aims to present definitive findings and validate the central hypotheses, the conclusion of this qualification document serves to consolidate the current progress, formalize the architectural decisions made during the design phase, and outline the necessary steps to complete the research.

\section{Current State of the Research}

To date, this research has successfully completed the foundational phases of the Design Science Research Methodology (DSRM). The problem of observability fragmentation in cloud-native environments was identified, and the theoretical requirements for achieving complete system visibility—grounded in control theory and distributed systems principles—were established.

The primary artifact of this thesis, the PANOPTES Reference Architecture, has transitioned from a conceptual blueprint to a concrete, operationalized testbed. The experimental environment, designed to avoid the multi-tenancy interference typical of public clouds, was successfully implemented using Infrastructure as Code (Vagrant and Ansible). Furthermore, the architecture itself was instantiated using Helm, integrating OpenTelemetry, Prometheus, Loki, Tempo, and Grafana into a cohesive telemetry pipeline monitoring a heterogeneous microservices workload (Online Boutique).

\section{Consolidated Architectural Decisions}

During the \textit{Design and Development} cycle, several critical architectural decisions were made and validated. These decisions, which will form part of the final contribution to the knowledge base, include:

\begin{itemize}
    \item \textbf{Adoption of K3s over Kubeadm:} To guarantee that the measured resource overhead is attributable to the observability stack and not to orchestration bloat, K3s was selected due to its minimal footprint, ensuring functional stability within the constraints of the virtualized laboratory.
    \item \textbf{Transition to Tempo for Tracing:} While Jaeger was initially considered, the architecture pivoted to Grafana Tempo. This decision was driven by Tempo's native semantic integration with Loki and Prometheus within the Grafana UI, fulfilling the requirement for a unified ``single pane of glass'' and mitigating the cognitive load of switching between disparate visualization tools.
    \item \textbf{Declarative Fault Injection:} The adoption of Chaos Mesh as a Kubernetes Operator ensures that the stress scenarios required for the controlled experiment are deterministic, versionable, and strictly reproducible.
\end{itemize}

\section{Remaining Challenges and Future Work}

The architecture is currently deployed and functional; however, the formal evaluation phase (DSRM Step 5) is pending execution. The immediate next step is the systematic execution of the controlled experiment defined in Chapter \ref{chap:methodology}. 

This involves injecting faults via Chaos Mesh and rigorously measuring the Mean Time to Detect (MTTD) and the resource overhead ($CPU/RAM$) across both the baseline (Vanilla Kubernetes) and the experimental (PANOPTES) groups. The statistical analysis of these results will dictate whether the formulated hypotheses ($H_1$ and $H_2$) are accepted or rejected.

\section{Schedule and Milestones}

To ensure the timely completion of this doctoral research leading up to the final defense, the following milestones have been established:

\begin{itemize}
    \item \textbf{Month 1-2 (Post-Qualification):} Systematic execution of the controlled experiments ($2\times2$ factorial design) and raw data collection.
    \item \textbf{Month 3:} Statistical analysis of the results (Hypothesis testing) and drafting of the Results and Discussion chapter.
    \item \textbf{Month 4:} Submission of the PANOPTES architecture and preliminary results to a peer-reviewed conference or journal (e.g., IEEE/ACM).
    \item \textbf{Month 5:} Final revision of the thesis document, incorporating the feedback from the qualification committee.
    \item \textbf{Month 6:} Final Thesis Defense.
\end{itemize}

By adhering to this structured plan, this research is positioned to contribute robust, empirical evidence to the field of cloud-native observability, providing the academic community and practitioners with a validated, open-source architectural standard.